\section{Los dos grafos}\label{sec:grafos}

El contador mira el contenido de lo que sale. Los grafos miran el contexto: qué tan sensible es lo
que la sesión está tocando, y si lo que sale desciende de algo que se leyó antes. Los dos están en
\cod{engine/graph/}.

Conviene decir de entrada cuál es su papel real, porque es más modesto que el que tenían en el diseño.
En la versión evaluada los grafos aportan \emph{señales por llamada} al modelo de evidencia
(\cref{sec:evidencia}). La parte que propaga sospecha entre pasos lejanos está implementada y
probada, pero ninguna regla de decisión la lee (\cref{sec:no-conectado}).

\subsection{El grafo de archivos}\label{sec:gff}

\paragraph{Qué representa.} Un grafo no dirigido $G_R$ del repositorio de trabajo
(\cod{engine/graph/gff.py}). Sus nodos son los archivos, las carpetas, las claves del manifiesto y los
destinos permitidos de cada clave. Tiene tres clases de aristas, cada una con un peso
(\cref{tab:aristas}).

\begin{table}[ht]
\caption{Las aristas del grafo de archivos.}\label{tab:aristas}
\small
\begin{tabularx}{\textwidth}{@{}P{3.4cm}LR{1.2cm}@{}}
\toprule
Clase & Une & Peso \\
\midrule
Árbol de carpetas & Cada archivo o carpeta con la carpeta que lo contiene & 1 \\
Importación & Un archivo Python con cada módulo del repositorio que importa (se lee con el analizador
  sintáctico de Python, sin ejecutar nada) & 2 \\
Referencia & Cada clave con los archivos donde vive, y con sus destinos permitidos & 4 \\
\bottomrule
\end{tabularx}
\end{table}

El grafo se construye una vez, al arrancar. Del manifiesto usa los identificadores, los archivos y los
destinos; el valor de las claves no se lee.

\paragraph{Qué calcula.} A cada nodo se le quiere asignar una \emph{sensibilidad} entre 0 y 1. El
manifiesto etiqueta unos pocos: las rutas sensibles (\cod{hot_paths}, con valor 1: los archivos de
claves) y las que se sabe que no lo son (\cod{cold_paths}, con valor 0: la documentación). El resto
se infiere suponiendo que nodos vecinos tienen sensibilidades parecidas. Es regresión con un campo
gaussiano sobre el grafo~\cite{zhu}: el campo tiene como matriz de precisión el laplaciano del grafo
más un término de regularización,
\begin{equation}
  Q = D - A + \kappa I ,
  \label{eq:Q}
\end{equation}
donde $A$ es la matriz de pesos, $D$ la diagonal de sus sumas por fila y $\kappa = \num{0.05}$. Si $P_B$
es la matriz diagonal que marca los nodos etiquetados, $y_B$ sus valores y $\tau = \num{0.2}$ el ruido de
las etiquetas, la distribución posterior es gaussiana con
\begin{equation}
  Q_{\mathrm{post}} = Q + \tau^{-2} P_B , \qquad
  Q_{\mathrm{post}}\, \mu = \tau^{-2} P_B\, y_B , \qquad
  \sigma^2 = \diag\!\left(Q_{\mathrm{post}}^{-1}\right) .
  \label{eq:post}
\end{equation}
La media $\mu$ es la sensibilidad estimada de cada nodo y $\sigma$ su incertidumbre. Se resuelve con
una factorización LU dispersa; la diagonal de la inversa se obtiene resolviendo contra columnas de la
identidad, en bloques de 256.

\paragraph{Qué entrega en cada llamada.} La función \cod{touch} recibe los recursos de la llamada y
devuelve tres números.
\begin{description}
  \item[$m$.] La sensibilidad $\mu$ del nodo más sensible que toca la llamada.
  \item[$m_{\mathrm{sd}}$.] La incertidumbre $\sigma$ de ese mismo nodo.
  \item[\cod{jump}.] La distancia, en saltos del grafo, entre lo que toca esta llamada y lo que
    tocaron las cinco anteriores. Mide si la sesión saltó a una zona del repositorio donde no estaba.
\end{description}
Si la llamada no toca ningún nodo del grafo, los tres valen 0. Es el caso de una petición a un
\esp{host} que no es destino permitido de ninguna clave, y resulta ser una señal informativa
(\cref{sec:modelo-ajustado}).

\paragraph{El grafo de la demo.} El repositorio de la demo es chico: 30 nodos (13 archivos, 7
carpetas, 5 claves y 5 destinos) y 30 aristas (19 de árbol, 1 de importación y 10 de referencia), con
5 nodos etiquetados. La \cref{tab:grafo} muestra el resultado en algunos de ellos.

\begin{table}[ht]
\caption{Sensibilidad posterior de algunos nodos del repositorio de la demo. Calculado con
\cod{FileGraph.from_workspace} sobre \cod{demo/repo}.}\label{tab:grafo}
\small
\begin{tabular}{@{}lllrr@{}}
\toprule
Nodo & Clase & Etiqueta del manifiesto & $\mu$ & $\sigma$ \\
\midrule
\tabladatos{tab_grafo}
\bottomrule
\end{tabular}
\end{table}

La tabla se lee así. Los archivos de claves quedan cerca de 1 y con poca incertidumbre, porque están
etiquetados. La clave y su destino permitido heredan una sensibilidad alta por las aristas de
referencia. La documentación queda cerca de 0. Una carpeta que contiene a la vez un archivo sensible
y uno que no lo es (\cod{config/}) queda en el medio. Y el código, que está lejos de toda etiqueta,
queda en un valor intermedio con mucha incertidumbre: el modelo dice que no sabe.

\subsection{El grafo de llamadas}\label{sec:dag}

\paragraph{Qué representa.} Un grafo dirigido acíclico de las llamadas de una sesión
(\cod{engine/graph/dag.py}). Hay una arista de la llamada $i$ a la llamada posterior $j$ cuando los
argumentos de $j$ comparten una subcadena de $k = 8$ o más caracteres con el resultado de $i$, en
claro o en una de sus variantes (base64 en sus tres alineaciones, hexadecimal, codificación de URL,
invertido y ROT13). La arista dice: algo de lo que devolvió $i$ viajó en lo que envió $j$.

\paragraph{Las dos etiquetas.} Cada llamada lleva dos etiquetas que se heredan por las aristas.
\begin{description}
  \item[$U$.] La llamada es no confiable o desciende de una que lo es.
  \item[$S$.] El conjunto de claves que la llamada leyó o que le llegaron por sus ancestros. Una
    llamada lee una clave si toca su archivo o si su resultado contiene 8 caracteres seguidos de su
    valor vigente.
\end{description}
La etiqueta $U$ viaja por cualquier arista. La etiqueta $S$ de una clave viaja por una arista solo si
el material compartido incluye 8 caracteres seguidos del valor vigente de esa clave. Sin esa
condición, $S$ se heredaba con cualquier texto en común y casi todo quedaba marcado.

\paragraph{El patrón que detecta.} Un sumidero que tiene $U$ y tiene $S$ no vacío: lo que sale
desciende a la vez de contenido no confiable y de la lectura de una clave. En ese caso el grafo
devuelve un \emph{camino}: la llamada no confiable más antigua, la lectura más antigua y el sumidero.
Es la única información del grafo de llamadas que usa una regla de decisión (la regla 3 de la
\cref{sec:decision}), y sirve también para explicar la decisión a una persona.

\paragraph{El índice.} Buscar subcadenas comunes contra todas las llamadas anteriores sería
cuadrático. El grafo mantiene un índice de los $k$-gramas de todos los resultados vistos: un
\esp{hash} polinomial de cada ventana de 8 caracteres apunta a la llamada dueña. El índice son bloques
ordenados que se fusionan por tamaño, de modo que consultar los $m$ $k$-gramas de unos argumentos
cuesta $O(m \log N)$ sin recorrer las llamadas anteriores. Toda coincidencia de \esp{hash} se confirma
con la subcadena real antes de crear la arista. Cada resultado se indexa hasta \num{200000}
caracteres. Según la bitácora del detector, este cambio bajó el percentil 95 del tiempo por llamada
de 15,6 a 0,55~ms en sesiones de \num{1000} llamadas.

\paragraph{Su límite.} Un envío en pedazos de menos de 8 caracteres no crea arista: ningún pedazo
comparte 8 caracteres con lo leído. El goteo, que es el caso que motiva el proyecto, es invisible para
este grafo. Lo cubren el contador y la rotación.

\subsection{Lo que está implementado y no se usa para decidir}\label{sec:no-conectado}

El diseño original tenía una pieza más ambiciosa, y está en el repositorio con sus tests.

\begin{description}
  \item[Aristas de sospecha y $R^{+}$ (\cod{suspicion.py}).] Además de las aristas de contenido, que
    son seguras, se agregan aristas \emph{probables} entre una fuente (una llamada no confiable o que
    leyó una clave) y cualquier llamada posterior. Cada una se abre con probabilidad
    $p_{ij} = g \cdot e^{-(t_j - t_i)/\tau_s}$, con $\tau_s = 50$ llamadas, donde $g$ depende de la
    sensibilidad de lo que tocan las dos llamadas y de su correlación posterior en el grafo de
    archivos. $R^{+}(c)$ es la probabilidad de que exista un camino desde contenido no confiable,
    pasando por la lectura de una clave, hasta la llamada $c$. Se calcula de forma exacta cuando hay
    hasta 16 aristas inciertas y por Monte Carlo (200 muestras) en otro caso.
  \item[Explicación por dominadores (\cod{explain.py}).] Para un sumidero bloqueado, calcula las
    llamadas por las que pasa \emph{todo} camino desde una fuente no confiable: la explicación mínima y
    el mejor punto de corte. Usa la intersección de Cooper, Harvey y Kennedy~\cite{chk} en una sola
    pasada, porque el orden temporal ya es un orden topológico.
\end{description}

\begin{limite}
\textbf{Ninguna regla lee $R^{+}$.} En el detector evaluado, el campo \cod{R_plus} de la decisión vale
1 si hay camino de contenido y 0 si no. Una investigación posterior, hecha solo sobre validación,
encontró que conectar $R^{+}$ no cambiaría nada medible con estos datos: no marca ningún envío que la
regla simple no marque (\cod{results/test/LECTURA.md}). Queda como trabajo futuro, y así lo decimos:
los grafos «conectan pasos lejanos» en el código, no en las decisiones.
\end{limite}
